\documentclass[11pt]{article}
\usepackage[margin=1in]{geometry}
\usepackage[T1]{fontenc}
\usepackage[utf8]{inputenc}
\usepackage{lmodern,microtype}
\usepackage{amsmath,amssymb,amsthm,mathtools}
\usepackage{booktabs,longtable,array,enumitem,graphicx,xcolor}
\usepackage{hyperref}
\hypersetup{colorlinks=true,linkcolor=blue,citecolor=blue,urlcolor=blue,pdftitle={Finite-State Theorems for Distance-Constrained Dynamic Subtraction Games}}
\setlength{\parindent}{0pt}\setlength{\parskip}{.5em}
\setlist[itemize]{leftmargin=2em}\setlist[enumerate]{leftmargin=2em}
\newtheorem{theorem}{Theorem}[section]
\newtheorem{lemma}[theorem]{Lemma}
\newtheorem{corollary}[theorem]{Corollary}
\newtheorem{proposition}[theorem]{Proposition}
\newtheorem{remark}[theorem]{Remark}
\newtheorem{conjecture}[theorem]{Conjecture}
\DeclareMathOperator{\mex}{mex}
\newcommand{\K}{\{1,\ldots,k\}}
\newcommand{\Pn}{\mathcal{P}}
\newcommand{\Nn}{\mathcal{N}}

\title{\textbf{Distance-Constrained Dynamic Subtraction Games:}\\
An Interval Structure Theorem and a Certified Affine Regime}
\author{Research Team\\\small Computational research manuscript --- author review required}
\date{9 October 2026 --- reproducible research draft}

\begin{document}
\maketitle
\begin{abstract}
We consider a one-heap normal-play subtraction game in which a player may remove
at most $k$ counters, but a removal $r$ following a removal $s$ must satisfy
$|r-s|>m$.  We show that the losing previous-move states at any fixed heap size
form an integer interval (possibly empty), which for $m=1$ has width at most
three unless it is the entire previous-move set.  We also establish eventual
periodicity of the full previous-move outcome sequence for every fixed pair
$(k,m)$ by a finite-memory argument.  These statements form Theorem A and
admit elementary mathematical proofs.  For the principal case $m=1$ we then
give a separately labeled, reproducible \emph{computer-assisted finite-range
proof} (Theorem B): for each integer $58\le k\le78$, the unrestricted starting
heap sizes that are losing are exactly
$n\equiv0,86,k+87\pmod{2k+88}$.  An independently implemented
Boolean interval recurrence, finite $k$-row cycle witnesses, a verified
finite-prefix bridge, and direct Sprague--Grundy cross-checks supply the
certificate.  This does not yet give a uniform symbolic explanation of the
affine period or the endpoints $58,78$.  We make this distinction explicit,
provide checkable evidence, and formulate the remaining proof problem.
\end{abstract}

\textbf{Keywords:} combinatorial game theory; impartial games; subtraction
games; P-positions; eventual periodicity; finite-state certificate.
\tableofcontents
\newpage

\section{Introduction and research question}
Ordinary subtraction games use a prescribed set of permissible removals.
Repetition-free subtraction games add a dependence on the previous removal;
Gurvich's study of that setting is a direct point of comparison
\cite{gurvich2012}.  We investigate the modified constraint that consecutive
removals must differ by more than a fixed nonnegative integer $m$.

A first computational investigation revealed structured periodic sets of
initial losing heap sizes as the maximum removal $k$ varies.  The objective
is twofold: to find rigorous general results about the states and to prove
specific parameter-dependent period formulas.  We distinguish four levels
of mathematical status throughout this manuscript:
\begin{itemize}
\item \textbf{Definition or deduction:} a statement proved from the game rules.
\item \textbf{Finite computation:} an output observed in a bounded calculation.
\item \textbf{Certified fixed-parameter result:} an infinite statement for a
fixed parameter established from a finite, independently checkable certificate
plus a proved recurrence invariant.
\item \textbf{Uniform symbolic theorem:} a result verified by a mathematical
argument valid for all parameters in the stated family without enumerating
the individual parameter values.
\end{itemize}

Theorem A belongs to the first category. Theorem B belongs to the third;
a uniform symbolic proof and a parameter-boundary explanation remain open.
This terminology also prevents a finite-tail period detector from being
mistaken for a periodicity proof.

\section{Game definition and normal-play outcomes}
Let $k\ge1$ and $m\ge0$ be integers.  In $D(k,m)$, a heap contains $n\ge0$
counters.  The current player removes an integer
$1\le r\le\min(k,n)$.  Except on the first turn, the previous removal $s$
is part of the state and legality requires $|r-s|>m$.
The player with no legal move loses; this includes cases with counters left
but all removal sizes forbidden.

After a first move a state is $(n,s)$ with $s\in\K$.  Its legal children are
\[
\mathcal{M}_{k,m}(n,s)
=\{(n-r,r): 1\le r\le\min(k,n),\;|r-s|>m\}.
\]
Let $P(n,s)$ mean that $(n,s)$ is a previous-player winning position,
equivalently losing for the player to move.  Then
\begin{equation}\label{eq:pn}
P(n,s)\iff\forall r\in\{1,\ldots,\min(k,n)\}
\text{ with }|r-s|>m,\quad \neg P(n-r,r).
\end{equation}
The empty universal quantification is true, so every $(0,s)$ is P.
For the unrestricted starting heap, write $P_0(n)$ for the same outcome
without a previously played removal.

For context, the standard Sprague--Grundy recurrence assigns
\begin{equation}\label{eq:grundy}
g(n,s)=\mex\{g(n-r,r):(n-r,r)\in\mathcal{M}_{k,m}(n,s)\},
\end{equation}
and $g_0(n)=\mex\{g(n-r,r):1\le r\le\min(k,n)\}$.
The classical theory implies $P(n,s)$ exactly when $g(n,s)=0$
\cite{sprague1935,grundy1939,winningways2001}.  Our Boolean solver determines
whether $g$ is zero; it does \emph{not} compute its nonzero values.

\section{Theorem A: exact row structure and eventual periodicity}
For fixed $k,m$ define the set of P-state previous removals
\[
Q_n=\{s\in\K:P(n,s)\}
\]
and the \emph{potential P-reaching moves}
\begin{equation}\label{eq:An}
A_n=\{r\in\{1,\ldots,\min(k,n)\}:r\in Q_{n-r}\}.
\end{equation}
The word ``potential'' matters: the membership test does not impose the
current previous-move constraint.  Rather, it records all removals reaching a
P-child, from which legality for each $s$ can be recovered.

\begin{lemma}[Boolean outcome invariant]\label{lem:invariant}
For every $n\ge0$, $s\in\K$ and $m\ge0$,
\begin{equation}\label{eq:invariant}
P(n,s)\iff A_n\subseteq\{s-m,\ldots,s+m\},
\qquad
P_0(n)\iff A_n=\varnothing.
\end{equation}
\end{lemma}
\begin{proof}
A move $r$ reaches a P-child exactly when $r\in A_n$.
From $(n,s)$ it is additionally legal exactly when $|r-s|>m$.
Thus $(n,s)$ is P if and only if no member of $A_n$ is legal,
which is precisely the first inclusion.  At the first turn no removal is
excluded by a previous move, so the unrestricted position is P precisely
when there is no P-reaching removal, i.e., $A_n=\varnothing$.
\end{proof}

\begin{theorem}[Theorem A1: interval row theorem]\label{thm:interval}
For all integers $k\ge1$, $m\ge0$ and $n\ge0$, $Q_n$ is either empty or an
integer interval inside $\K$.  More precisely, if $A_n=\varnothing$, then
$Q_n=\K$.  Otherwise, letting $a_n=\min A_n$ and $b_n=\max A_n$,
\begin{equation}\label{eq:interval}
\boxed{
Q_n=\bigl[\max(1,b_n-m),\;\min(k,a_n+m)\bigr]\cap\mathbb Z,
}
\end{equation}
where a lower endpoint greater than its upper endpoint denotes the empty
set.  If $A_n\ne\varnothing$ and $Q_n\ne\varnothing$, its cardinality is
at most $2m+1$.  In particular, for $m=1$ it has at most three elements.
\end{theorem}
\begin{proof}
By Lemma~\ref{lem:invariant}, each $r\in A_n$ must satisfy
$r-m\le s\le r+m$.  Intersecting these inequalities yields
\[
\max_{r\in A_n}(r-m)\le s\le\min_{r\in A_n}(r+m),
\]
which equals $b_n-m\le s\le a_n+m$.  Clipping to
$1\le s\le k$ proves the formula.  If nonempty, its width is bounded by
$(a_n+m)-(b_n-m)+1\le2m+1$ because $a_n\le b_n$.
\end{proof}

\begin{corollary}[Compact outcome alphabet]\label{cor:alphabet}
For each fixed $k,m$, the full P/N row $Q_n$ takes values among
$1+k(k+1)/2$ possibilities (the empty set and all nonempty integer
intervals).  For $m=1$, all non-full, nonempty rows contain at most
three successive previous-move values.
\end{corollary}

\begin{theorem}[Theorem A2: eventual periodicity]\label{thm:eventual}
For every fixed pair of integers $k\ge1$ and $m\ge0$, there exist
integers $t\ge0$ and $p\ge1$ such that
\[
\boxed{Q_{n+p}=Q_n\qquad(n\ge t).}
\]
Moreover, the unrestricted starting outcome $P_0(n)$ is eventually periodic
with a period dividing $p$.
\end{theorem}
\begin{proof}
For $n\ge k$, every candidate removal is in $\K$, and
Equation~\eqref{eq:An} constructs $A_n$ entirely from
$Q_{n-1},\ldots,Q_{n-k}$.  By Theorem~\ref{thm:interval},
$Q_n$ is then a deterministic function of those previous $k$ rows.
Consequently the window
$W_n=(Q_{n-k+1},\ldots,Q_n)$ evolves deterministically for $n\ge k-1$.
Each row is one of finitely many subsets of $\K$, so there are finitely
many windows.  The infinite sequence of windows must repeat; deterministic
evolution forces subsequent windows to repeat with the same positive shift.
This yields $Q_{n+p}=Q_n$ for all sufficiently large $n$.
Finally, $P_0(n)$ is equivalent to $A_n=\varnothing$ by
Lemma~\ref{lem:invariant}.  Once all rows $Q_{n-r}$ for $1\le r\le k$
lie in the periodic tail, the same computation gives
$A_{n+p}=A_n$ and hence $P_0(n+p)=P_0(n)$.
Its least period therefore divides the full-state period $p$.
\end{proof}

\begin{remark}
The proof establishes eventual periodicity but gives neither the least
period nor a useful description of its dependence on $k$.  Theorem B
addresses that harder quantitative question over one finite parameter
interval.  Theorem A is proved directly from the game's rules; standard
P/N and Grundy theory is cited for context and terminology.
\end{remark}

\section{Algorithms, verification methodology, and original data}
\subsection{Independent computational methods}
The research repository contains the original Python SG engine, an optimized
C++17 SG scanner, a Boolean bitset solver, an exact affine geometric
extractor, and an integer-interval row classifier.  The present audit adds
\texttt{src/verify\_theorems\_ab.py} to verify the two formal results
independently of the original period-scanning heuristic.

The interval recurrence uses Lemma~\ref{lem:invariant} and
Theorem~\ref{thm:interval}.  If $A_n$ is empty, write a full P row;
otherwise write the clipped interval in Equation~\eqref{eq:interval}.
The code computes every $A_n$ from previously computed rows and requires no
floating-point arithmetic.

An unrelated \emph{naive} Boolean implementation explicitly quantifies over
legal moves for each state.  The two implementations are compared for
$k=1,\ldots,9$, $m=0,1,2,3$, and $0\le n\le75$ (36 parameter cases).
For the 21 main parameter values, the interval rows and unrestricted
starting flags are also compared against direct SG evaluation over the
finite range used in the certificate.

\subsection{Why a long numerical scan is not a proof}
The exploratory C++ study scanned $m=1$, $35\le k\le120$ through
$n=500{,}000$.  It found, among other structures, an affine regime
\[
42\le k\le52:\quad p=2k+56,\quad
P_0(n)\iff n\equiv0,54,k+55\pmod p,
\]
and the regime targeted below.  These were discoveries, not certificates.
In particular, a shorter $k=80$ scan suggested a misleading long transient,
whereas a later deep scan produced a different period candidate.  The
underlying lesson is that an observed repeated tail does not itself prove
infinite continuation.

\section{Finite-window certificate lemma}
\begin{lemma}[Repeated windows force a full-state cycle]\label{lem:window}
Fix $k,m$ and let $p\ge1$.  Suppose for some $t\ge0$ that
\begin{equation}\label{eq:window}
Q_{t+i}=Q_{t+p+i}\quad\text{for every }0\le i<k.
\end{equation}
Then $Q_{n+p}=Q_n$ for every $n\ge t$.
\end{lemma}
\begin{proof}
Equation~\eqref{eq:window} supplies the first $k$ equal pairs.
For $n\ge t+k$, the rows needed to compute $Q_n$ are
$Q_{n-r}$, $1\le r\le k$.  Each has a matching row $Q_{n-r+p}$
by the inductive hypothesis, so the computed potential move sets
$A_n$ and $A_{n+p}$ are equal.  Theorem~\ref{thm:interval}
then gives $Q_n=Q_{n+p}$.  The first inductive step has $n\ge k$,
so the unrestricted removal range is already $1,\ldots,k$.
\end{proof}

\begin{lemma}[Finite bridge for unrestricted starting positions]
\label{lem:bridge}
Under the assumptions of Lemma~\ref{lem:window}, the initial P/N sequence
is $p$-periodic for all $n\ge t+k$.
If a proposed $p$-periodic formula agrees with $P_0(n)$ for every
\begin{equation}\label{eq:bridge}
0\le n\le t+k+p-1,
\end{equation}
then it agrees for all $n\ge0$.
\end{lemma}
\begin{proof}
For $n\ge t+k$, the candidate removals in $A_n$ are all $1\le r\le k$,
while their child rows satisfy $Q_{n+p-r}=Q_{n-r}$ by
Lemma~\ref{lem:window}.  Hence $A_{n+p}=A_n$ and
$P_0(n+p)=P_0(n)$.  The tested interval contains the entire
nonperiodic prefix $0,\ldots,t+k-1$ and one full period of the
eventual periodic region $t+k,\ldots,t+k+p-1$.  Repeated subtraction
of $p$ reduces any later index to that verified period.
\end{proof}

\section{Theorem B: certified affine starting-position regime}
\begin{theorem}[Theorem B: exact finite-range affine classification]
\label{thm:B}
For each integer $k$ with $58\le k\le78$, consider the unrestricted
starting state of $D(k,1)$ with $n\ge0$ counters.  Then this state is a
P-position if and only if
\begin{equation}\label{eq:B}
\boxed{n\equiv0,\ 86,\ \text{or }k+87\pmod{2k+88}.}
\end{equation}
The least positive period of the unrestricted binary P/N sequence is
$2k+88$.

\emph{Proof status: finite-range computer-assisted theorem.}  The statement
is infinite in $n$, but its 21 integer parameter instances are established
by the finite recurrence checks in the accompanying audit script and the
inductive Lemmas~\ref{lem:window}--\ref{lem:bridge}; no uniform symbolic
formula for all full-state rows is claimed.
\end{theorem}

\begin{proof}[Computer-assisted proof of Theorem B]
Fix one $k\in\{58,\ldots,78\}$ and set $p=2k+88$.  The script
\texttt{src/verify\_theorems\_ab.py} computes $Q_n$ using the exact
integer recurrence of Theorem~\ref{thm:interval}.  It searches for
$t\in\{0,1,\ldots,k\}$ satisfying the $k$ exact row equalities
\eqref{eq:window}, and asserts that a witness exists.
The computed $t$ for each $k$ is recorded in
Table~\ref{tab:cert}.  Lemma~\ref{lem:window} proves from those equalities
that the \emph{full-state} rows repeat forever beyond $t$.

Next the script compares every computed unrestricted P/N flag for
$0\le n\le t+k+p-1$ with the predicate
\[
(n\bmod p)\in\{0,86,k+87\}.
\]
All comparisons pass for each $k$ in the stated range.  By
Lemma~\ref{lem:bridge}, these finitely many checks imply
Equation~\eqref{eq:B} for every $n\ge0$.

For least-period minimality, let $R=\{0,86,k+87\}\subset\mathbb Z/p\mathbb Z$.
If the periodic binary sequence had a smaller period, the
three-element residue set would be invariant under some shift whose
orbit length is a proper divisor of three, necessarily an orbit of
length three.  Thus $R$ would have to be the equally spaced set
$\{0,p/3,2p/3\}$.  But throughout $58\le k\le78$,
\[
68\le p/3\le244/3<86< 2p/3,
\]
so $86$ cannot be one of the two nonzero equally spaced residues.
Therefore $p$ is the least positive period.

The results of all 21 finite computations are included in the generated
audit table, alongside reproducible program source and CSV hashes.  Direct
SG cross-checks independently agree on every computed outcome through the
required bridge bound.  This completes the stated finite-range,
computer-assisted proof, conditional only on the audit code executing as
provided; it is \emph{not} a parameter-uniform hand derivation of the
formula.
\end{proof}

\subsection{Complete fixed-parameter certificate table}
\input{generated/certificate_rows.tex}

\subsection{Representative full-state structure}
For $k=68$, the formula gives $p=224$ and starting P-residues
$\{0,86,155\}$.  The computed full-state diagram, rather than only the
unrestricted starting outcomes, exhibits full rows and short diagonal
interval families:
\begin{center}
\IfFileExists{../figures/full_state_k68_p_positions.png}{%
  \includegraphics[width=.95\linewidth]{../figures/full_state_k68_p_positions.png}
}{\emph{Place the repository's figures directory next to paper to render the heatmap.}}
\end{center}
The figure is exploratory evidence for a concise eventual symbolic
classification; it is not used as a proof of Theorem~\ref{thm:B}.

\section{Exact row mining and the remaining symbolic proof problem}
The geometric extraction initially identified exact line segments,
including structures described by $n-s=C$, but did not classify all
P-state cells using its selected families.  Applying
Theorem~\ref{thm:interval} row-by-row eliminates that geometric omission:
all computed P-states can be represented exactly as full rows, empty rows,
or short integer intervals.

The separate \texttt{row\_interval\_classifier.py} and its CSV outputs
compress many rows into affine blocks and track recurring formulas across
$k=58,\ldots,78$.  Such a representation is an exact \emph{computed}
classification of those finite state tables.  Some rows still require
explicit exceptions, and the compressed representation is not yet a
uniform nonrecursive formula for all $n$ and all permitted $k$.

A natural next target is to construct a piecewise function
\begin{equation}\label{eq:F}
F_k(u)\subseteq\K,\qquad 0\le u<2k+88,
\end{equation}
using finitely many affine interval formulas.  To prove it symbolically,
one must show that the associated P-reaching set
\[
B_n=\{r\in\{1,\ldots,\min(k,n)\}:r\in F_k((n-r)\bmod p)\}
\]
regenerates $F_k(n\bmod p)$ by Theorem~\ref{thm:interval}, with
special attention to the full-state transient for $k=58,\ldots,64$.
This is a well-defined collection of symbolic inequalities in $n,k$.
It is not solved by the certificates in this paper.

\section{Boundary checks, limitations, and originality}
The same candidate formula fails at the neighboring parameters
$k=57$ and $k=79$ in finite computations, providing counterexamples
to extending Theorem B unchanged.  These failures suggest that the
bounds $58$ and $78$ reflect changes in the full-state geometry, but
we have not derived those bounds from general symbolic inequalities.

The observed regime $42\le k\le52$ exhibits the related empirical pattern
$p=2k+56$ and residue set $\{0,54,k+55\}$.  Both patterns fit the template
\[
p=2k+A+2,\qquad R=\{0,A,k+A+1\},
\]
with $A=54$ or $86$.  No general infinite family in $A$ is claimed.
Moreover, it has not yet been established that the variant $D(k,m)$ and
its particular affine classifications are new relative to all prior
literature.  Gurvich \cite{gurvich2012} explicitly studies
repetition-free subtraction games, so comparisons must be checked
carefully before making a novelty claim.

The calculation is finite and fully reproducible, but independent
replication, code review, and team-authored exposition remain essential.
In particular, a computer-assisted proof and a closed-form symbolic proof
answer different mathematical questions.

\section{Conclusion and completed workflow}
This research progressed through the following verifiable stages:
\begin{enumerate}
\item Define the history-dependent move constraint and enlarge states to $(n,s)$.
\item Implement a direct SG recurrence and an optimized independent scanner.
\item Identify candidate affine periods from deep scans, then distinguish
finite-tail observations from proven infinite continuation.
\item Replace nonzero SG tracking with the Boolean potential-move set $A_n$.
\item Prove Theorem A1 (interval row structure) and Theorem A2 (eventual
full-state and starting-position periodicity) directly.
\item Exploit finite memory: one repeated pair of complete $k$-row
windows proves indefinitely repeated future full-state outcomes.
\item Check the finite-prefix bridge for every $k=58,\ldots,78$, with
independent direct SG verification, to obtain Theorem B.
\item Explore affine and row-interval structure to identify the outstanding
parameter-uniform symbolic classification problem.
\end{enumerate}

\textbf{Conclusion.}  The general eventual periodicity and interval
structure are now rigorous mathematical statements.  The affine period
$2k+88$ and residue set $\{0,86,k+87\}$ are exactly established for all
21 specified integer values of $k$ by a checkable finite-state proof.
What remains open is a compact parameter-uniform explanation of the
full-state intervals and of the regime boundaries, followed by a broader
classification across additional parameter ranges.

\paragraph{Authorship and computational assistance.}
This document is an AI-assisted research draft, not a record of an
independent human-written submission.  The student team should independently
review the proofs, validate all computations, verify the bibliography,
and disclose computational/AI assistance as required by its institution or
competition.  The original external sources are credited separately below;
no external proof text is reproduced.

\section*{Works cited}
\addcontentsline{toc}{section}{Works cited}
\input{references.tex}

\appendix
\section{Reproduction instructions and certificate contract}
From the repository root, run
\begin{verbatim}
python src/verify_theorems_ab.py \
 --out data/certificates/theorem_B_audit.csv \
 --summary data/certificates/theorem_B_audit.json
\end{verbatim}
The checker verifies 36 small $(k,m)$ cases against an independently
written Boolean game recursion; for each of 21 primary $k$ values, it
computes exact interval rows and a repeated $k$-row window at distance
$p=2k+88$, checks all starting P-flags through $t+k+p-1$, cross-checks
outcomes against a direct SG computation through the same bound, and
checks period minimality.  This is the reproducible finite certificate
underlying Theorem B.  See \texttt{PROOF\_AUDIT.md} for its invariants.

\section{Research-history preservation}
The original repository paper from the oldest commit (4 October 2026)
is archived unchanged as \texttt{paper/archive/main\_initial\_2026-10-04.tex}.
Its earlier table's ``prefix checked through'' values were shorter than
the bound actually used to bridge from full-state periodicity to starting
periodicity; they have been replaced by the mathematically sufficient
$t+k+p-1$ bound in this draft.  The certificate results themselves did
not change. The historical archive is provided for provenance, not as
a source to be cited as an independently published result.
\end{document}
